% Definition file for row def_3_19 -- "DBlockedWalks".
% Auto-generated stub by scaffold/solve_chapter.py at row start. The
% manager agent (or a worker dispatched by it) fills the body below with the
% Lean-faithful definition statement(s). Stays close to the lecture notes.
%
% This file is a sibling of `main.tex` inside `<subsection>/tex/`. The
% `subfiles` package lets it compile both standalone (resolving its
% preamble via `main.tex`) and as part of `main.tex` via `\subfile{...}`.

\documentclass[main]{subfiles}

\begin{document}
% ref: def_3_19
% title: DBlockedWalks
\def\rowref{def\_3\_19}
\phantomsection\label{def_3_19}

\begin{Def}[$d$-blocked walks]\label{def:d-blocking}
      Let $G=(J,V,E,L)$ be a CDMG and $C \ins J \cup V$ a subset of nodes and $\pi$ a walk in $G$:
      \[ \pi =\lp  v_0 \sus \cdots \sus v_n \rp.  \]
      \begin{enumerate}
          \item
              We say that the walk $\pi$ is \emph{$C$-$d$-blocked} or \emph{$d$-blocked by $C$} to emphasize the use of the bidirected edges\footnote{The ``d'' here stands for ``directional''. d-separation was first only used for DAGs (without bidirected edges).  For ADMGs it was then called ``m-separation'' in \cite{Ric03}. But since the notion of m-separation is arguably the natural extension of d-separation and to avoid introducing more definitions, we will just call it d-separation as well, which will not create any ambiguity.}
              if either:
      \begin{enumerate}[label=\roman*.)]
          \item $v_0 \in C$ or $v_n \in C$ or:
          \item there are two adjacent edges in $\pi$ of one of the following forms:
              \[\begin{array}{rccl}
           \text{ left chain: } & v_{k-1} \hut v_k \hus v_{k+1} & \text{ with } & v_k \in C,\\
           \text{ right chain: } & v_{k-1} \suh v_k \tuh v_{k+1} & \text{ with } & v_k \in C,\\
           \text{ fork: } & v_{k-1} \hut v_k \tuh v_{k+1} & \text{ with } & v_k \in C,\\
           \text{ collider: } & v_{k-1} \suh v_k \hus v_{k+1} & \text{ with } & v_k \notin \Anc^G(C).
              \end{array}\]
      \end{enumerate}
          \item We say that the walk $\pi$ is \emph{$C$-$d$-open} if it is not $C$-$d$-blocked.
      \end{enumerate}
\end{Def}

\end{document}
